\documentclass[10pt]{article}

\def\Type{LessonDJ}
\def\TypeNum{Workshop Week 10} %Lesson #
\def\TypeTitle{Reflecting on Optimization } %Lesson Title
\def\Semester{Fall 2021} %Not Used 
\def\Course{MATH 1190}%Course number and name
\def\Targets{  \target{A2} }

\usepackage{preambleDJ} 

%%%%%%%%%%%% BEGIN DOCUMENT %%%%%%%%%%%%%%%
%%%%%%%%%%%% BEGIN DOCUMENT %%%%%%%%%%%%%%%
%%%%%%%%%%%% BEGIN DOCUMENT %%%%%%%%%%%%%%%
%%%%%%%%%%%% BEGIN DOCUMENT %%%%%%%%%%%%%%%
%%%%%%%%%%%% BEGIN DOCUMENT %%%%%%%%%%%%%%%

\begin{document}
\pagestyle{otherpages}
\thispagestyle{firstpage}
\vs {.6}
\Large
You work in a company where the daily revenue of the company, $R$ (in dollars)  depends on the temperature, $T$ (in $^{\circ}$C ),  of the air conditioning in the office. Reflect on the process of maximizing $R(T)$ below


%\mpageT{.55}{.45}{
{\bf Conceptual Discussion}

\tasksA{2}{
%\task Explain how you could use the idea of rate of change to identify where Q might reach its best value.
\task Why is the derivative the right tool to find the temperature $T$ that gives the highest $R$?
\task Describe in your own words the steps you would use to find the temperature $T$ that gives the highest  $R$?
}


\graphicsC{qrcode.jpeg}{4}

\end{document}